\documentclass[11pt,a4paper]{article}

% --- 1. CORE PACKAGES ---
\usepackage[utf8]{inputenc}
\usepackage[T1]{fontenc}
\usepackage[ngerman]{babel} % German language support
\usepackage[margin=1in]{geometry}
\usepackage{lmodern}
\usepackage{microtype}
\usepackage{amsmath,amssymb}
\usepackage{graphicx}
\usepackage{booktabs}
\usepackage{array}
\usepackage{enumitem}
\usepackage{fancyhdr}
\usepackage{titlesec}
\usepackage{xcolor}
\usepackage{listings}
\usepackage{tikz}
\usetikzlibrary{shapes.geometric, arrows, positionings, calc, shadows}
\usepackage[many]{tcolorbox}
\usepackage[colorlinks=true, linkcolor=NavyBlue, urlcolor=SlateBlue, citecolor=EmeraldGreen]{hyperref}

% --- 2. COLOR PALETTE DEFINITIONS (Corporate Bank Theme) ---
\definecolor{NavyBlue}{HTML}{0F172A}      % Slate 900 (Primary Dark)
\definecolor{SlateBlue}{HTML}{1E293B}     % Slate 800 (Secondary Accent)
\definecolor{GoldAccent}{HTML}{EAB308}    % Gold 500 (Vindobona Theme)
\definecolor{EmeraldGreen}{HTML}{10B981}  % Emerald 500 (Income/Security)
\definecolor{RubyRed}{HTML}{EF4444}       % Red 500 (Expense/Blocked)
\definecolor{SkyBlue}{HTML}{0EA5E9}       % Sky 500 (Info/Vercel)
\definecolor{LightBG}{HTML}{F8FAFC}       % Slate 50
\definecolor{BorderGray}{HTML}{E2E8F0}    % Slate 200

% --- 3. PAGE STYLING & HEADERS ---
\pagestyle{fancy}
\fancyhf{}
\renewcommand{\headrulewidth}{0.5pt}
\renewcommand{\footrulewidth}{0.5pt}
\fancyhead[L]{\textcolor{NavyBlue}{\textbf{Vindobona Pro FinTech}}}
\fancyhead[R]{\textcolor{SlateBlue}{Architektur-Handbuch}}
\fancyfoot[L]{\textcolor{SlateBlue}{\small Wiener Edition}}
\fancyfoot[C]{\textcolor{SlateBlue}{\small Seite \thepage}}
\fancyfoot[R]{\textcolor{SlateBlue}{\small Vertrauliches Dokument}}

% ---  titling adjustments ---
\titleformat{\section}
  {\color{NavyBlue}\normalfont\Large\bfseries}
  {\thesection}{1em}{}[{\titlerule[0.8pt]}]
\titleformat{\subsection}
  {\color{SlateBlue}\normalfont\large\bfseries}
  {\thesubsection}{1em}{}
\titleformat{\subsubsection}
  {\color{NavyBlue}\normalfont\normalsize\bfseries}
  {\thesubsubsection}{1em}{}

% --- 4. LISTINGS STYLING (Syntax Highlighting) ---
\lstdefinelanguage{JavaScript}{
  keywords={break, case, catch, continue, debugger, default, delete, do, else, false, finally, for, function, if, in, instanceof, new, null, return, switch, this, throw, true, try, typeof, var, void, while, with, let, const, async, await, class, export, import, from, default},
  morecomment=[l]{//},
  morecomment=[s]{/*}{*/},
  morestring=[b]',
  morestring=[b]"
}

\lstset{
  language=JavaScript,
  backgroundcolor=\color{LightBG},
  basicstyle=\ttfamily\small\color{NavyBlue},
  keywordstyle=\color{SkyBlue}\bfseries,
  commentstyle=\color{gray}\itshape,
  stringstyle=\color{EmeraldGreen},
  identifierstyle=\color{NavyBlue},
  numbers=left,
  numberstyle=\tiny\color{gray},
  stepnumber=1,
  numbersep=8pt,
  showstringspaces=false,
  breaklines=true,
  frame=single,
  rulecolor=\color{BorderGray},
  tabsize=2,
  captionpos=b
}

% --- 5. CUSTOM CALLOUT BOXES ---
\newtcolorbox{ImportanceBox}[1]{
  colback=LightBG,
  colframe=GoldAccent,
  fonttitle=\bfseries\color{NavyBlue},
  title={⭐ Wichtiger Hinweis: #1},
  arc=6pt,
  boxrule=1.5pt,
  left=12pt,
  right=12pt,
  top=8pt,
  bottom=8pt
}

\newtcolorbox{ReasonBox}[1]{
  colback=LightBG,
  colframe=EmeraldGreen,
  fonttitle=\bfseries\color{white},
  coltitle=white,
  title={✔️ Begründung \& Nutzen: #1},
  arc=6pt,
  boxrule=1.5pt,
  left=12pt,
  right=12pt,
  top=8pt,
  bottom=8pt
}

\newtcolorbox{WarningBox}[1]{
  colback=LightBG,
  colframe=RubyRed,
  fonttitle=\bfseries\color{white},
  coltitle=white,
  title={⚠️ Sicherheits-Warnung: #1},
  arc=6pt,
  boxrule=1.5pt,
  left=12pt,
  right=12pt,
  top=8pt,
  bottom=8pt
}

\newtcolorbox{FeatureCard}[2]{
  colback=white,
  colframe=SlateBlue,
  fonttitle=\bfseries\color{white},
  coltitle=white,
  title={#1},
  subtitle={#2},
  arc=4pt,
  boxrule=1pt,
  left=8pt,
  right=8pt,
  top=6pt,
  bottom=6pt
}

% --- 6. DOCUMENT BEGIN ---
\begin{document}

% --- FRONT TITLE PAGE ---
\begin{titlepage}
    \centering
    \vspace*{2cm}
    {\huge\bfseries \color{NavyBlue} 🏛️ Vindobona Pro FinTech}\\[0.6cm]
    {\Large\bfseries \color{SlateBlue} Handbuch für die Full-Stack Finanzportal-Architektur}\\[0.3cm]
    \textcolor{GoldAccent}{\rule{12cm}{2pt}}\\[1.5cm]
    
    \begin{tcolorbox}[colback=LightBG, colframe=NavyBlue, arc=8pt, width=14cm]
        \centering
        \small\color{NavyBlue}
        Dieses technische Handbuch bietet eine detaillierte Analyse der gesamten Anwendungsarchitektur. Es dokumentiert jede Frontend-React-Komponente (\texttt{.tsx}), Stylesheet (\texttt{.css}), Backend-Route, Sicherheits-Middleware, Datenbankstruktur und Bereitstellungsschnittstelle.
    \end{tcolorbox}
    
    \vspace*{3cm}
    \begin{tabular}{r@{\hspace{0.5cm}}l}
        \textbf{Autor:} & Andy Rugero \\
        \textbf{Standort:} & Wien, Österreich \\
        \textbf{Ziel-Stack:} & React, Express, Docker, Kubernetes, Nginx, Azure \\
        \textbf{Sicherheitsstandard:} & PSD2-Bankenregulierung, OTP-2FA, TLS \\
    \end{tabular}
    
    \vfill
    {\small Dokument generiert für die Overleaf-Compiler-Integration \\ \today}
\end{titlepage}

\newpage
\tableofcontents
\newpage

% --- SECTION 1: ARCHITECTURE OVERVIEW ---
\section{🔌 1. Allgemeine Systemarchitektur}
Die Anwendung Vindobona Pro basiert auf einem entkoppelten Client-Server-Architekturmodell. Diese Trennung gewährleistet Systemleistung, Skalierbarkeit und klare Sicherheitsgrenzen.

\subsection{Client-Schicht (Frontend)}
Das Webportal verwendet einen Single-Page-Client (SPA), der mit \textbf{React, TypeScript und Vite} erstellt wurde.
\begin{itemize}
    \item \textbf{Typsicherheit:} TypeScript stellt sicher, dass komplexe Transaktionsdatenstrukturen zur Laufzeit mit den Payloads übereinstimmen und verhindert Abstürze.
    \item \textbf{Styling-Systeme:} Entwickelt mit \textbf{Vanilla CSS} und modularen CSS-Variablen. Vermeidet aufgeblähten Framework-Code, um eine maximale Render-Geschwindigkeit zu garantieren.
    \item \textbf{Zustandsverwaltung:} Der Anwendungszustand wird dynamisch über benutzerdefinierte React-Kontexte (\texttt{TransactionContext}) verwaltet, wodurch Ledger-Änderungen global über alle Anzeigemodule hinweg koordiniert werden.
\end{itemize}

\subsection{API-Schicht (Backend)}
Die Server-Engine basiert auf \textbf{Node.js und Express.js}.
\begin{itemize}
    \item \textbf{Authentifizierung:} Verwendet signierte, kurzlebige JSON Web Tokens (JWT), die sicher über Header oder HTTP-only Cookies übertragen werden.
    \item \textbf{Verifizierung:} Schützt Endpunkte vor Angriffen durch IP-basierte Ratenbegrenzungen (Rate Limiting), Brute-Force-Filter und strenge Rollenprüfungen.
\end{itemize}

\subsection{Datenbankschicht (Storage)}
Das System basiert auf einem \textbf{Dual-Dialekt-SQL-Treibermodell}:
\begin{itemize}
    \item \textbf{Entwicklung:} SQLite verarbeitet lokale Abfragen effizient in flachen Datenbankdateien.
    \item \textbf{Produktion:} Microsoft Azure hostet PostgreSQL-Datenbanken für die skalierbare Enterprise-Nutzung.
    \item \textbf{Abfrage-Übersetzung:} Eine dynamische Übersetzungs-Middleware fängt SQLite-Parameter (\texttt{?}) ab und tauscht sie zur Laufzeit in PostgreSQL-Äquivalente (\texttt{\$1, \$2}) aus.
\end{itemize}

\begin{ImportanceBox}{Entkoppelte Software-Architektur}
Die Trennung statischer Client-Dateien von der API-Verarbeitungsschicht ermöglicht effizientes Edge-Caching und reduziert die Last auf den Hauptservern.
\end{ImportanceBox}

\begin{ReasonBox}{Dual-Dialekt-Übersetzung}
Die Verwendung von lokalem SQLite beschleunigt Test-Setups und automatisierte Integrationsläufe. Die dynamische Syntax-Übersetzung in PostgreSQL ermöglicht ein nahtloses Deployment auf Produktionsdatenbanken in Azure, ohne dass Routen-Handler neu geschrieben werden müssen.
\end{ReasonBox}

\newpage

% --- SECTION 2: FRONTEND COMPONENT DIRECTORY ---
\section{🎨 2. Frontend-Komponenten \& Styling-Verzeichnis (\texttt{src/})}
Jedes Frontend-Element besteht aus einer TypeScript-Komponente und einem zugehörigen Stylesheet:

\subsection{Authentifizierung \& Sicherheits-Guards}
\begin{FeatureCard}{AuthScreen.tsx}{gepaart mit AuthScreen.css}
\textbf{Zweck:} Verwaltet Registrierung, sicheren Login, 2FA-Eingabe und Passwörter-Reset. \\
\textbf{Logik:} Kommuniziert mit \texttt{POST /api/users/register}, \texttt{POST /api/auth/login} und \texttt{POST /api/auth/2fa/login}. Bietet Formulare für Google Authenticator 2FA.
\end{FeatureCard}

\subsection{Layout-Komponenten}
\begin{FeatureCard}{Sidebar.tsx}{gepaart mit Sidebar.css}
\textbf{Zweck:} Zeigt das Benutzerprofil, das Avatar-Bild und das Navigationsmenü an. Schränkt den Zugriff auf das Admin-Panel basierend auf der Benutzerrolle ein.
\end{FeatureCard}

\begin{FeatureCard}{Topbar.tsx}{gepaart mit Topbar.css}
\textbf{Zweck:} Rendert Breadcrumbs, Theme-Wechsler (Hell/Dunkel), Logout-Button und das Mailbox-System für System-Updates.
\end{FeatureCard}

\subsection{Dashboard-Statistiken}
\begin{FeatureCard}{StatsRow.tsx \& StatCard.tsx}{gepaart mit StatsRow.css \& StatCard.css}
\textbf{Zweck:} Flex-Grid-Karten zur Visualisierung wichtiger Finanzkennzahlen (Kontostand, Einnahmen, Ausgaben) mit farblicher Statusanzeige.
\end{FeatureCard}

\subsection{Ledger \& Transaktionen}
\begin{FeatureCard}{TransactionList.tsx}{gepaart mit TransactionList.css}
\textbf{Zweck:} Tabelle zur Auflistung historischer Ledger-Einträge. Bietet CSV- und PDF-Exportfunktionen für Bankbelege.
\end{FeatureCard}

\begin{FeatureCard}{TransactionForm.tsx}{gepaart mit TransactionForm.css}
\textbf{Zweck:} Formular zur direkten Eingabe neuer Transaktionen mit Kategorie-Zuordnung.
\end{FeatureCard}

\subsection{Spezialisierte UI-Module}
\begin{FeatureCard}{CardWidget.tsx}{gepaart mit CardWidget.css}
\textbf{Zweck:} Ein interaktives 3D-Debitkarten-Widget, das Mausbewegungen folgt. Ermöglicht das sofortige Einfrieren der Karte.
\end{FeatureCard}

\begin{FeatureCard}{FXConverter.tsx}{gepaart mit FXConverter.css}
\textbf{Zweck:} Währungsumrechner für Transaktionen zwischen verschiedenen Währungskonten (EUR/USD/GBP) mit Live-Wechselkursen.
\end{FeatureCard}

\begin{FeatureCard}{Chatbot.tsx}{gepaart mit Chatbot.css}
\textbf{Zweck:} Schwebender AI-Finanzassistent zur Beantwortung von Fragen bezüglich Kontobewegungen per natürlicher Sprache.
\end{FeatureCard}

\newpage

% --- SECTION 3: BACKEND ARCHITECTURE ---
\section{⚙️ 3. Backend-Server- \& Service-Architektur (\texttt{backend/})}
Die Backend-Schicht ist nach dem Controller-Service-Muster strukturiert, um Transaktionssicherheit zu garantieren:

\subsection{Datenbank-Tabellen (Schema-DDL)}
Das Datenbankschema besteht aus fünf relationalen Haupttabellen:
\begin{lstlisting}[caption={Datenbank-Tabellendefinition}]
CREATE TABLE IF NOT EXISTS users (
    id TEXT PRIMARY KEY,
    username TEXT UNIQUE NOT NULL,
    email TEXT UNIQUE,
    password_hash TEXT NOT NULL,
    is_verified INTEGER DEFAULT 0,
    two_factor_secret TEXT,
    two_factor_enabled INTEGER DEFAULT 0,
    balance REAL DEFAULT 0.0,
    role TEXT DEFAULT 'user',
    is_card_frozen INTEGER DEFAULT 0,
    avatar_url TEXT
);
\end{lstlisting}

\subsection{Routen-Handler (Endpoints)}
\begin{itemize}
    \item \textbf{\texttt{auth.js}:} Steuert die Registrierung, Verifizierungscodes, Passwort-Resets und 2FA-Authentifizierung.
    \item \textbf{\texttt{transactions.js}:} Verarbeitet Ledger-Einträge, PDF-Downloads, Währungskonvertierungen und P2P-Transfers.
    \item \textbf{\texttt{admin.js}:} Administrator-Controller für Benutzer-Listen, Rollen-Upgrades und das Einsehen der Audit-Logs.
\end{itemize}

\subsection{Sicherheits-Middleware}
\begin{ImportanceBox}{authGuard.js}
Überprüft eingehende JWT-Tokens im Authorization-Header und injiziert die entschlüsselten Benutzerdaten in das Request-Objekt (\texttt{req.user}).
\end{ImportanceBox}

\begin{WarningBox}{PSD2-Transaktionssignierung}
Jeder ausgehende Geldtransfer erfordert zwingend die Eingabe eines gültigen 2FA-OTP-Codes vor der Ausführung der SQL-Transaktion (Strong Customer Authentication).
\end{WarningBox}

\newpage

% --- SECTION 4: CLOUD CONNECTIONS ---
\section{🚢 4. Produktions-Cloud-Topologie \& Netzwerkintegration}
Zur Gewährleistung von Hochverfügbarkeit und Ausfallsicherheit ist das System über eine verteilte Cloud-Architektur bereitgestellt:

\subsection{System-Komponenten-Mapping}
\begin{table}[h]
\centering
\caption{Deployment Systemtopologie}
\begin{tabular}{llll}
\toprule
\textbf{Schicht} & \textbf{Plattform} & \textbf{Rolle} & \textbf{Verbindungsweg} \\
\midrule
\color{SkyBlue}Frontend Client & Vercel CDN & Bereitstellung der UI & Browser $\rightarrow$ HTTPS (CDN) \\
\color{GoldAccent}API Server & Azure App Service & Ausführung des API-Codes & Client-JS $\rightarrow$ REST-API \\
\color{EmeraldGreen}Datenbank & Azure Database (Postgres) & Datenhaltung & API-Container $\rightarrow$ Port 5432 \\
\color{RubyRed}Reverse Proxy & Nginx / Kubernetes & Gateway \& Routing & Browser $\rightarrow$ Nginx $\rightarrow$ Docker \\
\bottomrule
\end{tabular}
\end{table}

\subsection{Netzwerk- und Integrationsflüsse}

\subsubsection{1. Source Code CI/CD-Pipelines (GitHub $\rightarrow$ Vercel \& Azure)}
\begin{itemize}
    \item \textbf{Vercel Frontend:} Commits auf dem \texttt{main}-Branch lösen Webhooks aus, die das React-Projekt bauen und weltweit auf Edge-Servern spiegeln.
    \item \textbf{Azure Backend:} GitHub Actions kompilieren das Backend, bauen ein \textbf{Docker-Image}, laden es auf Docker Hub hoch und stoßen das Azure App Service Update an.
\end{itemize}

\subsubsection{2. Laufzeitumgebungen (Die Rolle von Docker)}
\begin{ReasonBox}{Docker-Containerisierung}
Die Kapselung von Node.js-Laufzeitumgebungen, Bibliotheken und Code in einem Docker-Image stellt sicher, dass die API lokal, in CI/CD-Tests und in der Azure-Cloud identisch ausgeführt wird.
\end{ReasonBox}

\subsubsection{3. Ablauf eines Netzwerk-Requests}
\begin{enumerate}
    \item Der Browser lädt die statische React-Oberfläche verschlüsselt von \textbf{Vercel} (Port 443).
    \item Benutzeraktionen senden API-Anfragen an \texttt{api.vindobonafintech.com}.
    \item Die Anfrage trifft auf das **Nginx-Gateway**, welches das SSL-Zertifikat validiert und den Request intern an den **Docker-Container** auf **Azure** weiterleitet.
    \item Die Node-API im Docker-Container führt Abfragen auf der abgeschirmten **Azure PostgreSQL-Datenbank** aus und liefert JSON-Daten zurück.
\end{enumerate}

\newpage

% --- SECTION 5: DIAGRAMS ---
\section{🎨 5. Systemdiagramme \& Ablaufpläne}
Diese Sektion visualisiert die Anwendungsstrukturen mittels nativer LaTeX-TikZ-Diagramme:

\subsection{Allgemeine System-Architektur}
\begin{center}
\begin{tikzpicture}[node distance=1.8cm, auto, >=latex']
    % Styles
    \tikzstyle{client} = [rectangle, draw, fill=SkyBlue!20, text width=6cm, minimum height=1cm, rounded corners, align=center, draw=SkyBlue]
    \tikzstyle{server} = [rectangle, draw, fill=SlateBlue!20, text width=6cm, minimum height=1cm, rounded corners, align=center, draw=white, text=white]
    \tikzstyle{db} = [cylinder, draw, fill=EmeraldGreen!20, shape border rotate=90, text width=2cm, minimum height=1.5cm, align=center, draw=EmeraldGreen]
    \tikzstyle{ext} = [rectangle, draw, fill=GoldAccent!20, text width=3.5cm, minimum height=1cm, rounded corners, align=center, draw=GoldAccent]
    
    % Nodes
    \node (cl) [client] {\textbf{Client-Browser (Vercel CDNs)} \\ React, TSX, CSS};
    \node (ng) [server, below of=cl, yshift=-0.5cm] {\textbf{Nginx-Proxy / Kubernetes} \\ SSL-Terminierung \& Routing};
    \node (srv) [server, below of=ng, yshift=-0.5cm] {\textbf{Express-API (Azure/Docker)} \\ authGuard \& roleGuard Middleware};
    \node (db) [db, below of=srv, yshift=-0.8cm] {\textbf{Azure DB} \\ PostgreSQL};
    \node (email) [ext, right of=srv, xshift=4.5cm] {\textbf{SMTP Email} \\ Nodemailer};
    
    % Paths
    \draw[->, double, thick] (cl) -- (ng);
    \draw[->, thick] (ng) -- (srv);
    \draw[->, thick] (srv) -- (db);
    \draw[->, dashed, thick] (srv) -- (email);
\end{tikzpicture}
\end{center}

\subsection{Datenbank-Transaktionszyklus (Geldtransfer)}
\begin{center}
\begin{tikzpicture}[node distance=2.5cm, auto, >=latex']
    \tikzstyle{state} = [rectangle, draw, fill=LightBG, text width=7.5cm, minimum height=0.9cm, rounded corners, align=center, draw=NavyBlue]
    
    \node (s1) [state] {1. Client sendet Transfer-Request mit JWT \& 2FA-OTP};
    \node (s2) [state, below of=s1] {2. API validiert die Sitzung und prüft Kartenstatus};
    \node (s3) [state, below of=s2] {3. Server prüft Kontostand des Senders};
    \node (s4) [state, below of=s3] {4. BEGIN TRANSACTION startet atomaren DB-Lauf};
    \node (s5) [state, below of=s4] {5. Bucht Betrag um und schreibt Ledger-Einträge};
    \node (s6) [state, below of=s5] {6. COMMIT schreibt permanent $\rightarrow$ Bestätigung};
    
    \draw[->, thick] (s1) -- (s2);
    \draw[->, thick] (s2) -- (s3);
    \draw[->, thick] (s3) -- (s4);
    \draw[->, thick] (s4) -- (s5);
    \draw[->, thick] (s5) -- (s6);
\end{tikzpicture}
\end{center}

\end{document}
